\documentclass[10pt,a4paper]{amsart}
\usepackage{amsmath,amssymb,amsthm,graphicx}
% Retained preamble for the Stacks Project chapter snapshot in this directory.
% The chapter's own preamble.tex is a copy of the shared Stacks Project preamble; the block
% below is sliced MECHANICALLY, line for line, from it. Nothing was added, removed or
% redefined; the wrapper class, the xr-hyper/hyperref wiring and the externaldocument list
% are omitted because the assembler supplies the document class.
%
% lines taken from preamble.tex: 3, 5, 9, 11, 12, 156, 157, 158, ... (51 lines total)
%
% package coverage restored from the chapter's own preamble (2026-10-13): xy, xy-2cell

\usepackage[all]{xy}\usepackage{verbatim}
\xyoption{2cell}
\UseAllTwocells

\newenvironment{reference}{\comment}{\endcomment}
\newenvironment{slogan}{\comment}{\endcomment}
\newenvironment{history}{\comment}{\endcomment}
\theoremstyle{plain}
\newtheorem{theorem}[subsection]{Theorem}
\newtheorem{proposition}[subsection]{Proposition}
\newtheorem{lemma}[subsection]{Lemma}
\theoremstyle{definition}
\newtheorem{definition}[subsection]{Definition}
\newtheorem{example}[subsection]{Example}
\newtheorem{exercise}[subsection]{Exercise}
\newtheorem{situation}[subsection]{Situation}
\theoremstyle{remark}
\newtheorem{remark}[subsection]{Remark}
\newtheorem{remarks}[subsection]{Remarks}
\numberwithin{equation}{subsection}
\def\lim{\mathop{\mathrm{lim}}\nolimits}
\def\colim{\mathop{\mathrm{colim}}\nolimits}
\def\Spec{\mathop{\mathrm{Spec}}}
\def\Hom{\mathop{\mathrm{Hom}}\nolimits}
\def\Ext{\mathop{\mathrm{Ext}}\nolimits}
\def\SheafHom{\mathop{\mathcal{H}\!\mathit{om}}\nolimits}
\def\SheafExt{\mathop{\mathcal{E}\!\mathit{xt}}\nolimits}
\def\Sch{\mathit{Sch}}
\def\Mor{\mathop{\mathrm{Mor}}\nolimits}
\def\Ob{\mathop{\mathrm{Ob}}\nolimits}
\def\Sh{\mathop{\mathit{Sh}}\nolimits}
\def\NL{\mathop{N\!L}\nolimits}
\def\CH{\mathop{\mathrm{CH}}\nolimits}
\def\proetale{{pro\text{-}\acute{e}tale}}
\def\etale{{\acute{e}tale}}
\def\QCoh{\mathit{QCoh}}
\def\Ker{\mathop{\mathrm{Ker}}}
\def\Im{\mathop{\mathrm{Im}}}
\def\Coker{\mathop{\mathrm{Coker}}}
\def\Coim{\mathop{\mathrm{Coim}}}
\DeclareMathSymbol{\boxtimes}{\mathbin}{AMSa}{"02}
\def\QCohstack{\mathcal{QC}\!\mathit{oh}}
\def\Cohstack{\mathcal{C}\!\mathit{oh}}
\def\Spacesstack{\mathcal{S}\!\mathit{paces}}
\def\Quotfunctor{\mathrm{Quot}}
\def\Hilbfunctor{\mathrm{Hilb}}
\def\Curvesstack{\mathcal{C}\!\mathit{urves}}
\def\Polarizedstack{\mathcal{P}\!\mathit{olarized}}
\def\Complexesstack{\mathcal{C}\!\mathit{omplexes}}
\def\Pic{\mathop{\mathrm{Pic}}\nolimits}
\def\Picardstack{\mathcal{P}\!\mathit{ic}}
\def\Picardfunctor{\mathrm{Pic}}
\def\Deformationcategory{\mathcal{D}\!\mathit{ef}}
\usepackage{cleveref}
\crefname{originallemma}{original lemma}{original lemmas}
\Crefname{originallemma}{Original lemma}{Original lemmas}
\makeatletter
\newlabel{coherent-lemma-quasi-coherent-affine-cohomology-zero}{{01XB}{source}{}{sourcepointer-1}{}}
\newlabel{coherent-lemma-quasi-coherent-affine-cohomology-zero@cref}{{[originallemma][0][]01XB}{[1][0][]source}{}{}{}}
\makeatother
\crefname{originallemma}{original lemma}{original lemmas}
\Crefname{originallemma}{Original lemma}{Original lemmas}
\makeatletter
\newlabel{homology-lemma-characterize-weak-serre-subcategory}{{0754}{source}{}{sourcepointer-2}{}}
\newlabel{homology-lemma-characterize-weak-serre-subcategory@cref}{{[originallemma][0][]0754}{[1][0][]source}{}{}{}}
\makeatother
\crefname{originallemma}{original lemma}{original lemmas}
\Crefname{originallemma}{Original lemma}{Original lemmas}
\makeatletter
\newlabel{sites-cohomology-lemma-cartesian-flat-quasi-coherent}{{0GZR}{source}{}{sourcepointer-3}{}}
\newlabel{sites-cohomology-lemma-cartesian-flat-quasi-coherent@cref}{{[originallemma][0][]0GZR}{[1][0][]source}{}{}{}}
\makeatother
\crefname{originallemma}{original lemma}{original lemmas}
\Crefname{originallemma}{Original lemma}{Original lemmas}
\makeatletter
\newlabel{sites-cohomology-lemma-compare-topologies-derived-adequate-modules}{{07A8}{source}{}{sourcepointer-4}{}}
\newlabel{sites-cohomology-lemma-compare-topologies-derived-adequate-modules@cref}{{[originallemma][0][]07A8}{[1][0][]source}{}{}{}}
\makeatother
\crefname{originallemma}{original lemma}{original lemmas}
\Crefname{originallemma}{Original lemma}{Original lemmas}
\makeatletter
\newlabel{sites-modules-lemma-chaotic-quasi-coherent}{{0GZN}{source}{}{sourcepointer-5}{}}
\newlabel{sites-modules-lemma-chaotic-quasi-coherent@cref}{{[originallemma][0][]0GZN}{[1][0][]source}{}{}{}}
\makeatother
\crefname{originallemma}{original lemma}{original lemmas}
\Crefname{originallemma}{Original lemma}{Original lemmas}
\makeatletter
\newlabel{sites-modules-lemma-special-locally-free}{{03DM}{source}{}{sourcepointer-6}{}}
\newlabel{sites-modules-lemma-special-locally-free@cref}{{[originallemma][0][]03DM}{[1][0][]source}{}{}{}}
\makeatother
\crefname{originallemma}{original lemma}{original lemmas}
\Crefname{originallemma}{Original lemma}{Original lemmas}
\makeatletter
\newlabel{spaces-cohomology-section-higher-direct-image}{{071Y}{source}{}{sourcepointer-7}{}}
\newlabel{spaces-cohomology-section-higher-direct-image@cref}{{[originallemma][0][]071Y}{[1][0][]source}{}{}{}}
\makeatother
\crefname{originallemma}{original lemma}{original lemmas}
\Crefname{originallemma}{Original lemma}{Original lemmas}
\makeatletter
\newlabel{spaces-descent-lemma-quasi-coherent-alternative-small}{{0H04}{source}{}{sourcepointer-8}{}}
\newlabel{spaces-descent-lemma-quasi-coherent-alternative-small@cref}{{[originallemma][0][]0H04}{[1][0][]source}{}{}{}}
\makeatother
\crefname{originallemma}{original lemma}{original lemmas}
\Crefname{originallemma}{Original lemma}{Original lemmas}
\makeatletter
\newlabel{spaces-properties-lemma-properties-quasi-coherent}{{03M1}{source}{}{sourcepointer-9}{}}
\newlabel{spaces-properties-lemma-properties-quasi-coherent@cref}{{[originallemma][0][]03M1}{[1][0][]source}{}{}{}}
\makeatother
\title{Stacks, Derived Categories of Spaces: DQCoh-alternative-small}
\author{The Stacks Project Authors}
\date{stacks-project:spaces-perfect.tex:local-snapshot-20261012; GNU FDL 1.2 or later}
\begin{document}
\maketitle
\noindent\textit{Editorial note.} Statement span is lines 7596-7606 of sources/stacks-spaces-perfect/pinned-original-source.tex; proof span is lines 7608-7678, the complete author proof (69 body lines).\par\medskip
\noindent\textit{Editorial note.} Navigation stubs: 9 cited-but-not-included labels are supplied as non-load-bearing pointers with their pinned tags where the snapshot tag index has them.\par\medskip
\noindent\textit{Editorial note.} extra\_packages is empty on purpose: the retained preamble already loads xy with options and a second declaration clashes in the assembler wrapper.\par\medskip
\noindent\textit{Editorial note.} Text is retained verbatim, including the author's own typos and placeholders.\par\medskip
\noindent\textit{Editorial note.} Queue index 55 of queue\_spaceperfect.json (pb=3509, stmt 7596-7606, proof 7608-7678).\par\medskip
\noindent\textit{Editorial note.} stacks-spaces-perfect wave sp01 (2026-10-14): fresh adjudication accept per lane-s3/sp01\_judgments; pool per Chief 12:4x unfreeze (reserve-pool workflow, six-check clean); spaces-perfect.tex 126 with-proof rows, 70 b107-occupied labels avoided, 56 free rows judged 32 accept / 24 reject.\par\medskip
\noindent\textit{Editorial note.} Cross-chapter same-label warning (NOT an occupancy conflict, Chief ruling 2026-10-13 evening): label 'lemma-DQCoh-alternative-small' also exists in stacks-perfect — a different proposition in a different chapter file.\par\medskip
\noindent\textit{Allowed standard prerequisites (editorial).}\begin{itemize}
\item Explicit prerequisite boundary. The cited inputs are prerequisites; the argument in the retained span is not replaced by any cited result.
\end{itemize}\medskip
\ifdefined\hypertarget\hypertarget{sourcepointer-1}{}\fi
\noindent\textit{Source reference.} Original lemma 01XB: sources/stacks-spaces-perfect/pinned-original-source.tex, source lines 7678--7678. Cited by the retained statement or proof but not part of the collected spans; supplied only as a navigation pointer. The argument does not depend on its content.\par\medskip
\ifdefined\hypertarget\hypertarget{sourcepointer-2}{}\fi
\noindent\textit{Source reference.} Original lemma 0754: sources/stacks-spaces-perfect/pinned-original-source.tex, source lines 7678--7678. Cited by the retained statement or proof but not part of the collected spans; supplied only as a navigation pointer. The argument does not depend on its content.\par\medskip
\ifdefined\hypertarget\hypertarget{sourcepointer-3}{}\fi
\noindent\textit{Source reference.} Original lemma 0GZR: sources/stacks-spaces-perfect/pinned-original-source.tex, source lines 7678--7678. Cited by the retained statement or proof but not part of the collected spans; supplied only as a navigation pointer. The argument does not depend on its content.\par\medskip
\ifdefined\hypertarget\hypertarget{sourcepointer-4}{}\fi
\noindent\textit{Source reference.} Original lemma 07A8: sources/stacks-spaces-perfect/pinned-original-source.tex, source lines 7678--7678. Cited by the retained statement or proof but not part of the collected spans; supplied only as a navigation pointer. The argument does not depend on its content.\par\medskip
\ifdefined\hypertarget\hypertarget{sourcepointer-5}{}\fi
\noindent\textit{Source reference.} Original lemma 0GZN: sources/stacks-spaces-perfect/pinned-original-source.tex, source lines 7678--7678. Cited by the retained statement or proof but not part of the collected spans; supplied only as a navigation pointer. The argument does not depend on its content.\par\medskip
\ifdefined\hypertarget\hypertarget{sourcepointer-6}{}\fi
\noindent\textit{Source reference.} Original lemma 03DM: sources/stacks-spaces-perfect/pinned-original-source.tex, source lines 7678--7678. Cited by the retained statement or proof but not part of the collected spans; supplied only as a navigation pointer. The argument does not depend on its content.\par\medskip
\ifdefined\hypertarget\hypertarget{sourcepointer-7}{}\fi
\noindent\textit{Source reference.} Original lemma 071Y: sources/stacks-spaces-perfect/pinned-original-source.tex, source lines 7678--7678. Cited by the retained statement or proof but not part of the collected spans; supplied only as a navigation pointer. The argument does not depend on its content.\par\medskip
\ifdefined\hypertarget\hypertarget{sourcepointer-8}{}\fi
\noindent\textit{Source reference.} Original lemma 0H04: sources/stacks-spaces-perfect/pinned-original-source.tex, source lines 7678--7678. Cited by the retained statement or proof but not part of the collected spans; supplied only as a navigation pointer. The argument does not depend on its content.\par\medskip
\ifdefined\hypertarget\hypertarget{sourcepointer-9}{}\fi
\noindent\textit{Source reference.} Original lemma 03M1: sources/stacks-spaces-perfect/pinned-original-source.tex, source lines 7678--7678. Cited by the retained statement or proof but not part of the collected spans; supplied only as a navigation pointer. The argument does not depend on its content.\par\medskip
\begin{lemma}
\label{lemma-DQCoh-alternative-small}
In the situation above there are canonical exact equivalences between
the following triangulated categories
\begin{enumerate}
\item $D_\QCoh(\mathcal{O}_X)$,
\item $D_\QCoh(X_{affine, \etale}, \mathcal{O}_X)$,
\item $D_\QCoh(X_{affine}, \mathcal{O}_X)$, and
\item $\mathit{QC}(X_{affine}, \mathcal{O}_X)$.
\end{enumerate}
\end{lemma}

\begin{proof}
If $U \to V \to X$ are \'etale morphisms with $U$ and $V$ affine,
then the ring map $\mathcal{O}_X(V) \to \mathcal{O}_X(U)$ is flat. Hence
the equivalence between (3) and (4) is a special case of
Cohomology on Sites, Lemma
\ref{sites-cohomology-lemma-cartesian-flat-quasi-coherent}
(the proof also clarifies the statement).

\medskip\noindent
The discussion preceding the lemma shows that we have an equivalence of
ringed topoi
$(\Sh(X_{affine, \etale}), \mathcal{O}_X) \to (\Sh(X_\etale), \mathcal{O}_X)$
and hence an equivalence between abelian categories of modules.
Since the notion of quasi-coherent modules is intrinsic
(Modules on Sites, Lemma \ref{sites-modules-lemma-special-locally-free})
we see that this equivalence preserves the subcategories
of quasi-coherent modules. Thus  we get a canonical exact equivalence
between the triangulated categories in (1) and (2).

\medskip\noindent
To get an exact equivalence between the triangulated categories
in (2) and (3) we will apply Cohomology on Sites, Lemma
\ref{sites-cohomology-lemma-compare-topologies-derived-adequate-modules}
to the morphism $\epsilon : (X_{affine, \etale}, \mathcal{O}_X) \to
(X_{affine}, \mathcal{O}_X)$ above. We take
$\mathcal{B} = \Ob(X_{affine})$ and we take
$\mathcal{A} \subset \textit{PMod}(X_{affine}, \mathcal{O})$
to be the full subcategory of those presheaves $\mathcal{F}$
such that $\mathcal{F}(V) \otimes_{\mathcal{O}_X(V)} \mathcal{O}_X(U)
\to \mathcal{F}(U)$ is an isomorphism. Observe that by
Descent on Spaces, Lemma
\ref{spaces-descent-lemma-quasi-coherent-alternative-small}
objects of $\mathcal{A}$ are exactly those sheaves in the \'etale
topology which are quasi-coherent modules on
$(X_{affine, \etale}, \mathcal{O}_X)$.
On the other hand, by Modules on Sites, Lemma
\ref{sites-modules-lemma-chaotic-quasi-coherent},
the objects of $\mathcal{A}$ are exactly the quasi-coherent modules
on $(X_{affine}, \mathcal{O}_X)$, i.e., in the chaotic topology.
Thus if we show that Cohomology on Sites, Lemma
\ref{sites-cohomology-lemma-compare-topologies-derived-adequate-modules}
applies, then we do indeed get the canonical equivalence between the
categories of (2) and (3) using $\epsilon^*$ and $R\epsilon_*$.

\medskip\noindent
We have to verify 4 conditions:
\begin{enumerate}
\item Every object of $\mathcal{A}$ is a sheaf for the \'etale topology.
This we have seen above.
\item $\mathcal{A}$ is a weak Serre subcategory of
$\textit{Mod}(X_{affine, \etale}, \mathcal{O}_X)$. Above we have
seen that $\mathcal{A} = \QCoh(X_{affine, \etale}, \mathcal{O}_X)$ and
we have seen above that these,
via the equivalence $\textit{Mod}(X_{affine, \etale}, \mathcal{O}) =
\textit{Mod}(X_\etale, \mathcal{O}_X)$, correspond to the quasi-coherent
modules on $X$. Thus the result by Properties of Spaces, Lemma
\ref{spaces-properties-lemma-properties-quasi-coherent} and
Homology, Lemma \ref{homology-lemma-characterize-weak-serre-subcategory}.
\item Every object of $X_{affine}$ has a covering in the chaotic
topology whose members are elements of $\mathcal{B}$. This holds
because $\mathcal{B}$ contains all objects.
\item For every object $U$ of $X_{affine}$ and $\mathcal{F}$ in $\mathcal{A}$
we have $H^p_{Zar}(U, \mathcal{F}) = 0$ for $p > 0$.
This holds by the vanishing of cohomology of quasi-coherent modules
on affines, see discussion in Cohomology of Spaces, Section
\ref{spaces-cohomology-section-higher-direct-image} and
Cohomology of Schemes, Lemma
\ref{coherent-lemma-quasi-coherent-affine-cohomology-zero}.
\end{enumerate}
This finishes the proof.
\end{proof}

\end{document}
